\begin{longtable}[c]{ |c|c|c|c|c| }
\caption{根据 Key\_A、Key\_B 及 Key\_C 生成的 Key 值}
\label{tab:extmemencr-key-possibility}
\\\hline 
\rowcolor{lightgray}
Block$_A$ & Block$_B$ & Block$_C$ & Key  & Key 长度（位） \\\hline
\endfirsthead
\hline
\rowcolor{lightgray}
Block$_A$ & Block$_B$ & Block$_C$ & Key  & Key 长度（位） \\\hline
\endhead
\endfoot
存在 &   存在   &  无关项    &   \begin{math}\mathit{Key}_{A} || \mathit{Key}_{B}\end{math}    &   512   \\\hline
存在 &   不存在    &  无关项    &   \begin{math}\mathit{Key}_{A} || 0^{256}\end{math}    &   512   \\\hline
不存在  &   存在   &  无关项    &   \begin{math}0^{256} || \mathit{Key}_{B}\end{math}    &   512   \\\hline
不存在  &   不存在    &  存在       &   $\mathit{Key}_{C}$    &   256   \\\hline
不存在  &   不存在    &  不存在        &   $0^{256}$    &   256   \\\hline
\end{longtable}
\vspace{-1.5em}
{\bfseries 说明：}\\
\begin{itemize}
    \item “存在” 表示存在该 Block
    \item “不存在” 表示不存在该 Block
    \item ``$0^{256}$" 表示由 256 个位 0 组成的位串
    \item ``\begin{math}||\end{math}" 表示键合运算符，用于将一个位串连接到另一个位串
\end{itemize}
